\chapter{Аналитический раздел}


\section{Предметная область}
Кластеризация или кластерный анализ~\cite{Кластеризация} --- это задача разбиения множества объектов на группы, называемые кластерами. При этом внутри каждой группы должны оказаться похожие объекты, а объекты разных групп должны быть как можно более отличны. Основным отличием кластеризации от классификации является то, что перечень групп заранее не задан и определяется в процессе работы алгоритма.

Данные, с которыми работают методы кластеризации, по виду можно подразделить на числовые и категориальные.

Числовые данные~\cite{Данные} --- это данные, характеризующие состояние какого-либо параметра изучаемого объекта. Наиболее часто такие данные бывают представлены вещественными числами. Примерами числовых данных являются цена товара, вес человека, атмосферное давление, скорость автомобиля.

Категориальные данные~\cite{Данные} --- это данные, образующие признак принадлежности к какой-либо группе. Примерами категориальных данных являются цвет автомобиля, уровень образования человека.


\section{Меры расстояний}
Одним из важнейших аспектов кластеризации данных является выбор меры расстояний~\cite{Метрики}, которая определяет степень сходства или различия двух объектов.

Мерой расстояний между двумя объектами $d(x, y)$ называется неотрицательная функция, удовлетворяющая следующим правилам:
\begin{equation}
	\begin{cases}
		d(x, y) \ge 0 \\
		d(x, y) = d(y, x) \\
		d(x, y) = 0 \iff x = y \\
		d(x, y) \le d(x, u) + d(u, y)
	\end{cases},
\end{equation} 
где $x$, $y$, $u$ --- любые объекты кластерного анализа.

Мера расстояний может существенно повлиять на конечный результат кластеризации. Выбор меры расстояний в кластеризации должен осуществляться в зависимости от решаемой задачи и от типа данных. 

\subsection{Евклидово расстояние}
Евклидово расстояние~\cite{Кластеризация} --- наиболее широко используемая мера расстояний в кластеризации. Представляет собой расстояние между двумя точками в многомерном пространстве. Формула для евклидова расстояния:
\begin{equation}
	d(x, y) = \sqrt{\sum_{i=1}^m{{(x_i - y_i)}^2}},
\end{equation}
где $x_i$, $y_i$ --- координаты точек $x$, $y$ соответственно в $m$-мерном пространстве.

Несмотря на свою простоту, евклидово расстояние имеет ряд недостатков. Оно не инвариантно к масштабу, что вынуждает проводить нормализацию перед его использованием. Кроме того, с увеличением размерности пространства различия между расстояниями сглаживаются, что снижает информативность евклидовой меры~\cite{Евклид_размерность}. Применимо только для числовых данных.

\subsection{Квадрат евклидова расстояния}
Квадрат евклидова~\cite{Кластеризация} расстояния используется для придания большего веса более отдаленным друг от друга объектам. Вычисляется данное расстояние по следующей формуле:
\begin{equation}
	d(x, y) = \sum_{i=1}^m{{(x_i - y_i)}^2},
\end{equation}
где $x_i$, $y_i$ --- координаты точек $x$, $y$ соответственно в $m$-мерном пространстве.

Квадрат евклидова расстояния усиливает вклад удаленных объектов, что делает алгоритм более чувствительным к выбросам и при большом различии между объектами эта мера может привести к потери точности. Также как и евклидово расстояние квадрат евклидова расстояния подходит только для числовых данных.

\subsection{Расстояние городских кварталов}
Расстояние городских кварталов~\cite{Кластеризация} (манхэттенское расстояние) --- это сумма абсолютных разностей между декартовыми координатами точек. Высчитывается по формуле:
\begin{equation}
	d(x, y) = \sum_{i=1}^m{|x_i - y_i|},
\end{equation}
где $x_i$, $y_i$ --- координаты точек $x$, $y$ соответственно в $m$-мерном пространстве.

Название получено из-за особенности перемещения по городским кварталам, в которых единственными доступными направлениями являются вертикальное и горизонтальное. 

Расстояние городских кварталов часто демонстрирует более устойчивое поведение в пространствах высокой размерности, чем евклидово расстояние~\cite{Манхэттен_Евклид}. Для работы с категориальными данными данная мера не подходит.

\subsection{Расстояние Минковского}
Расстояние Минковского~\cite{distances} --- это обобщенная форма евклидова и манхэттенского расстояний. Расстояние Минковского порядка $p$ вычисляется по формуле:
\begin{equation}
	d(x, y) = (\sum_{i=1}^m{|x_i - y_i|^p})^{\frac{1}{p}}, p \ge 1,
	\label{eq::minc}
\end{equation}
где $x_i$, $y_i$ --- координаты точек $x$, $y$ соответственно в $m$-мерном пространстве, $p$ --- параметр, определяющий взаимодействие между признаками.

При $p = 1$ формула~\ref{eq::minc} принимает вид расстояния городских кварталов, при $p = 2$ --- расстояния Евклида.

Расстояние Минковского наследует недостатки метрик, которые оно обобщает. Наличие параметра $p$ позволяет гибко управлять поведением меры и выбирать оптимальное значение для конкретной задачи, однако подбор подходящего параметра может оказаться вычислительно затратным и требовать дополнительного анализа данных.

\subsection{Степенное расстояние}
Степенное расстояние~\cite{Кластеризация} является обобщением расстояния Минковского. Степенное расстояние вычисляется по формуле:
\begin{equation}
	d(x, y) = (\sum_{i=1}^m{|x_i - y_i|^p})^{\frac{1}{r}}, p \ge 1, r \ge 1,
	\label{eq::pow}
\end{equation}
где $x_i$, $y_i$ --- координаты точек $x$, $y$ соответственно в $m$-мерном пространстве.

Параметр $p$ в уравнении~\ref{eq::pow} отвечает за взвешивание разностей по отдельным координатам, параметр $r$ --- за взвешивание разностей между объектами. При $p = r$ мера принимает вид расстояния Минковского.

Наличие двух независимых параметров позволяет настраивать меру под разные задачи. Однако использование степенного расстояния требует тщательного подбора параметров $p$ и $r$, что может быть вычислительно затратно. Также как и расстояние Минковского степенное расстояние применимо только к числовым данным.

\subsection{Расстояние Чебышева}
Расстояние Чебышева~\cite{Кластеризация} определяется как наибольшая разность между значениями координат двух точек по любой координатной оси. Для нахождения расстояния Чебышева используется следующая формула:
\begin{equation}
	d(x, y) = \max_{i=1}^m{|x_i - y_i|},
\end{equation}
где $x_i$, $y_i$ --- координаты точек $x$, $y$ соответственно в $m$-мерном пространстве.

Расстояние Чебышева описывает минимальное число ходов, необходимых королю на шахматной доске для перехода из одной клетки в другую. Благодаря этому мера находит применение в задачах, где перемещение по различным направлениям считается равнозначным, например в некоторых игровых алгоритмах или моделях складской логистики. Расстояние Чебышева применимо только к числовым данным.

\subsection{Расстояние Хэмминга}
Расстояние Хэмминга~\cite{Хэмминг} --- это число позиций, различающихся между двумя векторами. Обычно его используют для сравнения двух двоичных строк одинаковой длины. Для двух векторов $\vec{x}$, $\vec{y}$ размерности $m$ расстояние Хэмминга вычисляется по формуле:
\begin{equation}
	d(\vec{x}, \vec{y}) = \sum_{i=1}^m{\delta(x_i, y_i)},
\end{equation}
где 
\begin{equation}
\delta(x_i, y_i) =
\begin{cases}
	0, & \text{если } x_i = y_i, \\
	1, & \text{если } x_i \neq y_i.
\end{cases}
\end{equation}

Расстояние Хэмминга учитывает лишь факт равенства или неравенства двух параметров, а не величину их различия. Данная мера часто применяется при выявлении ошибок в передаче данных и анализе бинарных последовательностей. Хотя расстояние Хэмминга можно применять как к числовым, так и к категориальным данным, для числовых признаков предпочтительно использовать другие меры расстояния.

\subsection{Расстояние Жаккара}
Расстояние Жаккара~\cite{Жаккар} используется для измерения степени различия между двумя множествами. Формула для вычисления расстояния Жаккара имеет вид:
\begin{equation}
	d(x, y) = 1 - \frac{|x \cap y|}{|x \cup y|},
\end{equation}
где $|x \cap y|$ — количество общих признаков между объектами $x$ и $y$, а $|x \cup y|$ — количество признаков, присутствующих хотя бы в одном из них.

Таким образом, расстояние Жаккара показывает, насколько объекты различаются относительно общего числа признаков. Чем больше доля совпадающих признаков, тем меньше значение расстояния.

мера Жаккара подходит для категориальных данных, однако она существенно зависит от размерности данных: при большом числе признаков схожие объекты могут демонстрировать заниженное сходство.


\subsection{Расстояние Гауэра}
Расстояние Гауэра~\cite{Гауэр} --- обобщенная мера, применимая для сравнения объектов как с числовыми признаками, так и c категориальными. Общий вид формулы расстояния Гауэра:
\begin{equation}
d(x, y) = \frac{\sum\limits_{i=1}^{m} ( s(x_i, y_i) \cdot \delta_i )}{\sum\limits_{i=1}^m \delta_i},
\end{equation}
где $m$ --- количество признаков, а $s(x_i, y_i)$ --- расстояние между $i$-м признаком объектов $x$ и $y$, $\delta_i$ принимает значение 1, если объекты $x$ и $y$ могут быть сравнены по $i$-ому признаку, и 0 иначе.

Если объекты $x$ и $y$ могут быть сравнены по всем признакам, то формула принимает вид:
\begin{equation}
	d(x, y) = \frac{1}{m} \sum\limits_{i=1}^{m} s(x_i, y_i),
\end{equation}

Для числовых признаков расстояние вычисляется по формуле:
\begin{equation}
	s(x_i, y_i) = \frac{|x_i - y_i|}{R_i},
\end{equation}
где $R_i$ --- диапазон значений $i$-ого признака.

Для категориальных признаков расстояние вычисляется по формуле:
\begin{equation}
	s(x_i, y_i) = 
	\begin{cases}
	0, & \text{если } x_i = y_i, \\
	1, & \text{иначе}.
	\end{cases}
\end{equation}

Расстояние Гауэра эффективно при работе с объектами, содержащими разные типы признаков. Его основным преимуществом является универсальность, так как мера позволяет обрабатывать как числовые, так и категориальные данные. Однако при большом количестве категориальных признаков вклад числовых признаков может снижаться из-за нормализации, что приводит к уменьшению чувствительности меры к числовым различиям.

\subsection{Сравнение мер расстояний}
Выделим следующие критерии для сравнения мер расстояний:
\begin{itemize}
	\item[---] К1 --- возможность находить расстояние между числовыми данными;
	\item[---] К2 --- возможность находить расстояние между категориальными данными.
\end{itemize}


В таблице~\ref{tbl::dst_meas} представлено сравнение рассмотренных мер расстояний по выделенным критериям.
\begin{table}[H]
	\begin{center}
		\caption{Сравнение мер расстояний}
		\begin{tabular}{|l|c|c|}
			\hline
			\textbf{Мера расстояний} & \textbf{К1} & \textbf{К2} \\
			\hline
			Евклидово расстояние & + & --- \\
			\hline
			Квадрат евклидова расстояния & + & --- \\
			\hline
			Манхэттенское расстояние & + & --- \\
			\hline
			Расстояние Минковского & + & --- \\
			\hline
			Степенное расстояние & + & --- \\
			\hline
			Расстояние Чебышева & + & --- \\
			\hline
			Расстояние Хэмминга & --- & + \\
			\hline
			Расстояние Жаккара & --- & + \\
			\hline
			Расстояние Гауэра & + & + \\
			\hline
		\end{tabular}
		\label{tbl::dst_meas}
	\end{center}
\end{table}

\section{Методы кластеризации}

\subsection{K-средних (k-means)}
Пусть задано конечное множество объектов $X = \{x_1, x_2, \dots, x_n\}$, которое необходимо разбить на $k$ кластеров $C = \{C_1, C_2, \dots, C_k\}$. Алгоритм k-средних~\cite{k_means} стремится минимизировать суммарное квадратичное отклонение точек кластеров от центров этих кластеров (центроидов):
\begin{equation}
	J_k = \sum_{i=1}^{k} \sum_{x \in C_i} ||x - \mu_i||^2,
	\label{eq::j_k}
\end{equation}
где $\mu_i$ --- центроид (среднее значение) кластера $C_i$, $||\cdot||$ --- евклидово расстояние.

Алгоритм состоит из следующих этапов:
\begin{enumerate}
	\item Случайным образом выбираются $k$ начальных центроидов.
	\item Каждый объект $x$ относится к кластеру, центроид которого находится ближе всего к нему.
	\item Пересчитываются координаты центроидов как среднее арифметическое всех объектов, попавших в соответствующий кластер.
	\item Шаги 2 и 3 повторяются до тех пор, пока центроиды не перестанут изменяться или изменение функции $J_k$ не станет меньше заданного порога~$\varepsilon$.
\end{enumerate}

Основными преимуществами метода являются линейная сложность $O(n)$ и сходимость при любых входных параметрах, что позволяет использовать его на больших выборках. Недостатками алгоритма являются необходимость заранее задавать число кластеров $k$, чувствительность к выбору начальных центроидов, возможность работы только с кластерами сферической формы.

\subsection{K-мод (k-modes)}
Алгоритм k-средних не применим к категориальным данным, так как для них невозможно найти среднее арифметическое. Алгоритм k-мод~\cite{k_modes} адаптирует идею k-средних для работы с категориальными данными.

Вместо евклидова расстояния в алгоритме обычно используется расстояние Хэмминга, а вместо вычисления среднего арифметического алгоритм определяет моду --- объект с наиболее часто встречающимися критериями. 

Общий вид алгоритма:
\begin{enumerate}
	\item Случайно выбираются $k$ объектов в качестве начальных мод.
	\item Каждый объект относится к кластеру, мода которого наиболее близка к нему. Для оценки сходства между объектами используется расстояние Хэмминга.
	\item Для каждого кластера вычисляется наиболее часто встречающиеся значения признаков. На их основе формируется новое значение моды.
	\item Шаги 2 и 3 повторяются до тех пор, пока мода не перестанет меняться.
\end{enumerate}

Алгоритм k-мод наследует преимущества и недостатки метода k-средних: он обладает линейной вычислительной сложностью, но требует заранее задавать число кластеров $k$, а также чувствителен к выбору начальных мод, что может приводить к нахождению локального, а не глобального оптимального решения. Также данный алгоритм способен получать кластеры только сферической формы.

\subsection{K-прототипов (k-prototypes)}
Алгоритм k-прототипов~\cite{k_prototypes} является обобщением алгоритмов k-средних и k-мод и позволяет работать с объектами, содержащими как числовые признаки, так и категориальные.

Для вычисления расстояния между объектами используется мера расстояний, которая является сочетанием квадрата евклидового расстояния для числовых признаков и расстояния Хэмминга для категориальных.

Общий вид алгоритма:
\begin{enumerate}
	\item Случайно выбираются $k$ объектов в качестве начальных центров.
	\item Каждый объект относится к кластеру, центр которого наиболее близок к нему.
	\item Вычисляется новый центр путем нахождения среднего значения для числовых признаков и моды для категориальных между объектами, принадлежащих кластеру.
	\item Шаги 2 и 3 повторяются до тех пор, пока центры не перестанут изменяться или алгоритм не достигнет максимального числа итераций.
\end{enumerate}

Алгоритм k-прототипов имеет те же недостатки, что k-средних и k-мод. Для него необходимо указывать число кластеров $k$, он чувствителен к выбору начальных центров и может получать кластеры только сферической формы. Однако данный алгоритм способен работать с данными, содержащими как числовые признаки, так и категориальные, что позволяет применять его для широкого круга задач кластеризации. Также он может обрабатывать большие объемы данных благодаря своей линейной сложности.

\subsection{K-медоидов (k-medoids)}
В отличие от k-средних, где центр кластера может быть фиктивной точкой пространства, в алгоритме k-медоидов~\cite{k_medoids} центром кластера всегда выбирается один из реальных объектов выборки --- медоид.

Медоид $\tilde{\mu_j}$ кластера $C_j$ выбирается таким образом, чтобы минимизировать сумму расстояний до остальных точек кластера:
\begin{equation}
	\tilde{\mu_j} = \operatorname*{argmin}_{y \in C_j} \sum_{x \in C_j} d(x, y),
\end{equation}
где $d(x, y)$ --- это выбранная мера расстояния между $x$ и $y$.

В общем случае нахождение нужного медоида является NP-полной задачей для точного решения~\cite{k_medoids_solve}. Поэтому на практике используются эвристические алгоритмы, которые находят приближенное решение. Наиболее распространенным подходом является алгоритм PAM (Partitioning Around Medoids).

Алгоритм PAM работает следующим образом:
\begin{enumerate}
	\item Случайно выбираются $k$ объектов в качестве медоидов $\{\tilde{\mu_1}, \dots, \tilde{\mu_k}\}$.
	\item Объекты распределяются по кластерам к ближайшему медоиду.
	\item Для каждого кластера выбирается новый медоид $\tilde{\mu_j}$ такой, что сумма $\sum\limits_{x \in C_j} d(x, \tilde{\mu_j})$ минимальна.
	\item Шаги 2 и 3 повторяются до тех пор, пока медоиды не перестанут изменяться.
\end{enumerate}

Алгоритм k-медоидов более устойчив к шуму и выбросам, чем k-средних, так как в качестве центра кластера рассматривается реальный объект, а не среднее арифметическое, которое чувствительно к выбивающимся данным. Однако вычислительная сложность алгоритма составляет $O(k(n-k)^2)$.

\subsection{C-средних (c-means)}
Алгоритм c-средних~\cite{c_means} относится к алгоритмам нечеткой кластеризации. В отличие от алгоритмов жесткой кластеризации, в которых каждый объект данных может принадлежать только одному кластеру, в алгоритмах нечеткой кластеризации объекты потенциально могут принадлежать нескольким кластерам.

Каждому объекту присваивается степень принадлежности. Эти степени принадлежности указывают на то, насколько объекты данных принадлежат каждому из кластеров. Объекты, находящиеся на окраине кластера и имеющие более низкую степень, принадлежат кластеру в меньшей степени, чем точки в центре кластера.

Алгоритм c-средних разбивает конечное множество объектов $X = \{x_1, \dots, x_n\}$ в набор из $k$ нечетких кластеров. Как и алгоритм k-средних, c-средних минимизирует целевую функцию:
\begin{equation}
	J_c = \sum_{i=1}^{n} \sum_{j = 1}^{k} w_{ij}^{m} ||x_i - v_j||^2,
	\label{eq::j_c}
\end{equation}
где:
\begin{itemize}
	\item[---] $v_j$ --- центроид $j$-ого кластера;
	\item[---] $w_{ij}$ --- степень принадлежности $i$-ого объекта $j$-ому кластеру, $w_{ij} \in [0, 1]$;
	\item[---] $m$ --- параметр нечеткости ($m > 1$).
\end{itemize}

Степень принадлежности вычисляется по формуле:
\begin{equation}
	w_{ij} = \frac{1}{\sum\limits_{r=1}^{k} \left( \frac{\|x_i - v_j\|}{\|x_i - v_r\|} \right)^{\frac{2}{m - 1}}}.
	\label{eq::w_i}
\end{equation}

Для расчета центроидов кластеров используется среднее взвешенное значение точек данных:
\begin{equation}
	v_{j} = \frac{\sum\limits_{i=1}^{n}w_{ij}^m x_i}{\sum\limits_{i=1}^{n} w_{ij}^m}.
	\label{eq::centre}
\end{equation}

Параметр $m$ отвечает за нечеткость кластеров: чем больше значение $m$, тем более размытые получаются кластеры. При $m \to 1$ степени принадлежности $w_{ij}$ сходятся к 0 или 1, а функция~\ref{eq::j_c} совпадает с функцией~\ref{eq::j_k} алгоритма k-средних. В отсутствии экспериментов или знаний в предметной области $m$ обычно принимают равным 2.

Общий вид алгоритма c-средних следующий:
\begin{enumerate}
	\item Случайно инициализируются степени принадлежности $w_{ij}$.
	\item Вычисляются центроиды кластеров по формуле~\ref{eq::centre}.
	\item Вычисляются степени принадлежности по формуле~\ref{eq::w_i}.
	\item Шаги 2 и 3 повторяются до тех пор, пока центры кластеров не перестанут изменяться или изменение целевой функции $J_c$ не станет меньше заданного порога $\epsilon$.
\end{enumerate}

В результате работы алгоритма будут получены множество центров кластеров $V = \{ v_1, \dots, v_k\}$ и матрица принадлежностей $W = [w_{ij}],\ i = 1, \dots, n,\ j = 1, \dots, k$.

Нечеткое разбиение данных позволяет получить более детальную информацию о принадлежности объектов кластерам, что широко применимо в биологии, химии, маркетинге и т.д., где объекты не принадлежат строго одному кластеру.

Основными плюсами алгоритма c-средних являются постоянная сходимость, возможность настройки нечеткости с помощью параметра $m$, более высокая устойчивость к выбросам, чем у алгоритма k-средних, за счет того, что центроиды вычисляются как взвешенные средние по степеням принадлежности.

Недостатками алгоритма являются зависимость от начальной инициализации степеней принадлежности, наличие параметра $m$, который требует настройки или поиска оптимального значения, увеличенное потребление памяти для хранения матрицы принадлежностей.


\subsection{DBSCAN}
DBSCAN~\cite{DBSCAN} (Density-Based Spatial Clustering of Applications with Noise) --- это алгоритм кластеризации на основе плотности, который группирует точки данных, расположенные близко друг к другу. Он идентифицирует кластеры как плотные области в пространстве данных, разделенные областями с меньшей плотностью.

DBSCAN использует два основных параметра: $\epsilon$ --- максимальное расстояние между двумя точками, при котором они считаются соседними, $minPts$ --- минимальное количество точек, необходимое для образования плотной области.

Алгоритм подразделяет точки на три основных класса:
\begin{itemize}
	\item[---] ядерные точки --- точки, которые имеют не менее $minPts$ других точек в пределах заданного расстояния $\epsilon$;
	\item[---] граничные точки --- точки, находящиеся в пределах заданного расстояния от ядерной точки, но не имеющие собственных $minPts$ соседей;
	\item[---] шумовые точки --- точки, которые не являются ни ядерными, ни граничными. Они находятся недостаточно близко к какому-либо кластеру, чтобы быть включенными в него.
\end{itemize}
Алгоритм DBSCAN выглядит следующим образом:
\begin{enumerate}
	\item Задаются параметры $\epsilon$ и $minPts$.
	\item Выбирается произвольная непосещенная точка из набора данных.
	\item Если точка имеет менее $minPts$ соседей в своей $\epsilon$-окрестности, то она помечается как шум (но может быть перемаркирована позже), иначе точка помечается как ядерная и инициализируется новый кластер.
	\item Все точки из $\epsilon$-окрестности ядерной точки добавляются в список обхода.
	\item Пока список обхода не пуст:
	\begin{enumerate}
		\item[5.1] Из списка обхода извлекается точка и, если она еще не была посещена, помечается как посещенная и включается в текущий кластер.
		\item[5.2] Если точка имеет не менее $minPts$ соседей в своей $\epsilon$-окрестности, то она является ядерной, и все ее соседи, которые еще не были добавлены в список обхода, добавляются в него.
		\item[5.3] Если точка не является ядерной, то она считается граничной и включается в кластер без дальнейшего расширения.
	\end{enumerate}
	\item Шаги 2–5 повторяются до тех пор, пока в наборе данных остаются непосещенные точки.
\end{enumerate}
Основными преимуществами DBSCAN являются способность находить кластеры произвольной формы без необходимости указания их количества заранее (в отличие от k-средних). Также алгоритм помечает выбросы как шум, что не искажает форму финальных кластеров.

К недостаткам относятся чувствительность к параметрам $\epsilon$ и $minPts$ (неправильный выбор может привести к слиянию или дроблению кластеров), неспособность справляться с кластерами сильно разной плотности и квадратичная вычислительная сложность $O(n^2)$ в худшем случае.


\subsection{Иерархическая кластеризация}

Иерархическая кластеризация~\cite{Иерархическая_кластеризация} --- это класс методов кластерного анализа, в которых структура кластеров представляется в виде иерархии, обычно изображаемой с помощью дендрограммы. В отличие от методов разбиения, иерархические алгоритмы не требуют предварительного задания числа кластеров.

Выделяют два класса методов иерархической кластеризации:
\begin{itemize}
	\item[---] агломеративная кластеризация --- метод, при котором изначально каждый объект образует отдельный кластер, после чего на каждом шаге происходит объединение двух наиболее близких кластеров;
	\item[---] дивизивная кластеризация --- метод, при котором изначально все объекты принадлежат одному единственному кластеру, который последовательно разделяется на более мелкие кластеры.
\end{itemize}

На практике чаще используется агломеративный подход из-за его более простой реализации.

На каждом шаге иерархической кластеризации необходимо вычислять расстояние между кластерами. Расстояние между одноэлементными кластерами определяется как расстояние между объектами этих кластеров. Для вычисления расстояния между кластерами $C_i$ и $C_j$ используются различные формулы в зависимости от спецификации. Наиболее часто используются следующие методы определения расстояния:

\begin{enumerate}
	\item Метод одиночной связи. Расстояние между двумя кластерами полагается равным минимальному расстоянию между двумя элементами из этих кластеров: 
	\begin{equation} 
		R_{min}(C_i, C_j) = min\{d(a, b) : a \in C_i, b \in C_j\},
	\end{equation}
	где $d(a, b)$ --- расстояние между элементами $a$ и $b$.
	\item Метод полной связи. Расстояние между двумя кластерами полагается равным максимальному расстоянию между двумя элементами из этих кластеров: 
	\begin{equation} 
		R_{max}(C_i, C_j) = max\{d(a, b) : a \in C_i, b \in C_j\}.
	\end{equation}
	\item Метод средней связи. Расстояние между двумя кластерами полагается равным среднему расстоянию между элементами этих кластеров:
	\begin{equation} 
		R_{avg}(C_i, C_j) = \frac{1}{|C_i| \cdot |C_j|} \sum_{a \in C_i} \sum_{b \in C_j} d(a, b),
	\end{equation}
	где $|C_i|$, $|C_j|$ --- число элементов в кластерах $C_i$ и $C_j$ соответственно.
	\item Центроидный метод. Расстояние между кластерами полагается равным расстоянию между их центроидами.
	\item Метод Уорда. Расстояние между кластерами полагается равным приросту суммы квадратов расстояний от точек до центроидов кластеров при объединении этих кластеров и вычисляется по формуле~\ref{eq::ward}. Данный метод стремится минимизировать внутрикластерную дисперсию.
	\begin{equation} 
		R_{ward}(C_i, C_j) = \frac{|C_i| \cdot |C_j|}{|C_i| + |C_j|} \cdot d^2(\mu_i, \mu_j),
		\label{eq::ward}
	\end{equation}
	где $\mu_i$, $\mu_j$ --- центроиды кластеров $C_i$, $C_j$ соответственно.
\end{enumerate}

Результатом работы алгоритма является дендрограмма, разрез которой на заданном уровне позволяет получить требуемое число кластеров.

Основными преимуществами иерархической кластеризации являются отсутствие необходимости заранее задавать число кластеров и возможность наглядной демонстрации результата. Однако при увеличении числа кластеров визуальная демонстрация теряет свою наглядность. Также к недостаткам относятся высокая вычислительная сложность $O(n^2)$ по памяти и времени и низкая устойчивость к шуму и выбросам.

\subsection{Смешанная гауссовская модель (GMM)}
Смешанная гауссовская модель~\cite{GMM} относится к методам нежесткой кластеризации. Предполагается, что наблюдаемые данные порождены совокупностью нескольких многомерных нормальных распределений.

Пусть задано конечное множество объектов
\[
X = \{x_1, x_2, \dots, x_n\}, \quad x_i \in \mathbb{R}^m,
\]
где $n$ --- число объектов, $m$ --- размерность пространства признаков. Требуется аппроксимировать распределение данных с использованием $K$ гауссовских компонент.

Смешанная гауссовская модель задается плотностью распределения
\begin{equation}
p(x) = \sum_{k=1}^{K} \pi_k \, \mathcal{N}(x \mid \mu_k, \Sigma_k),
\end{equation}
где:
\begin{itemize}
\item[---] $K$ — число гауссовских компонент;
\item[---] $\pi_k$ --- вес $k$-й компоненты, $\pi_k \ge 0$, $\sum\limits_{k=1}^{K} \pi_k = 1$;
\item[---] $\mu_k$ --- среднее значение $k$-й компоненты;
\item[---] $\Sigma_k$ --- ковариационная матрица $k$-й компоненты;
\item[---] $\mathcal{N}(x \mid \mu_k, \Sigma_k)$ — плотность многомерного нормального распределения.
\end{itemize}

Плотность многомерного нормального распределения имеет вид
\begin{equation}
\mathcal{N}(x \mid \mu_k, \Sigma_k) =
\frac{1}{(2\pi)^{m/2} |\Sigma_k|^{1/2}}
\exp\!\left(
-\frac{1}{2}(x - \mu_k)^\top \Sigma_k^{-1}(x - \mu_k)
\right).
\end{equation}

Для каждого объекта $x_i$ вводится скрытая переменная
\[
z_i \in \{1, 2, \dots, K\},
\]
которая указывает номер гауссовской компоненты, с которой связан объект $x_i$.

Вероятность принадлежности объекта $x_i$ $k$-й компоненте определяется по формуле Байеса:
\begin{equation}
P(z_i = k \mid x_i) =
\frac{\pi_k \, \mathcal{N}(x_i \mid \mu_k, \Sigma_k)}
{\sum\limits_{j=1}^{K} \pi_j \, \mathcal{N}(x_i \mid \mu_j, \Sigma_j)}.
\end{equation}

Параметры модели $\{\pi_k, \mu_k, \Sigma_k\}_{k=1}^{K}$ оцениваются методом максимального правдоподобия с использованием алгоритма ожидания-максимизации (EM).

Алгоритм EM включает следующие этапы:
\begin{enumerate}
\item Задаются начальные значения параметров $\pi_k$, $\mu_k$, $\Sigma_k$.
\item (E-шаг) Для каждого объекта $x_i$ вычисляются вероятности $P(z_i = k \mid x_i)$.
\item (M-шаг) Параметры модели обновляются по формулам:
\begin{equation}
\pi_k = \frac{1}{n} \sum_{i=1}^{n} P(z_i = k \mid x_i),
\end{equation}
\begin{equation}
\mu_k = \frac{\sum\limits_{i=1}^{n} P(z_i = k \mid x_i)\, x_i}
{\sum\limits_{i=1}^{n} P(z_i = k \mid x_i)},
\end{equation}
\begin{equation}
\Sigma_k =
\frac{\sum\limits_{i=1}^{n} P(z_i = k \mid x_i)\,
(x_i - \mu_k)(x_i - \mu_k)^\top}
{\sum\limits_{i=1}^{n} P(z_i = k \mid x_i)}.
\end{equation}
\item Шаги 2 и 3 повторяются до тех пор, пока изменение логарифма правдоподобия не станет меньше заданного порога $\varepsilon$ или не будет достигнуто максимальное число итераций.
\end{enumerate}

Результатом работы алгоритма является набор параметров смешанной гауссовской модели и матрица вероятностей принадлежности объектов гауссовским компонентам.

Смешанная гауссовская модель обладает рядом характерных свойств, определяемых предположениями, лежащими в основе модели. Поскольку каждая компонента описывается многомерным нормальным распределением, модель корректно аппроксимирует данные в случае, когда кластеры имеют эллипсоидальную форму, а плотность данных внутри кластеров изменяется плавно. В рамках кластеризации модель реализует нечеткий подход, при котором каждому объекту сопоставляется вероятность принадлежности к компонентам, что допускает перекрывающиеся кластеры и неопределенные границы между ними.

В то же время применение смешанной гауссовской модели связано с рядом ограничений. Модель предполагает гауссовскую природу распределения данных внутри каждой компоненты, поэтому кластеры  неэллиптической формы, а также данные с выраженными областями разной плотности могут быть аппроксимированы некорректно. Кроме того, модель ориентирована на работу с непрерывными числовыми признаками и не подходит для категориальных данных без соответствующей модификации распределений. Результат кластеризации также зависит от выбора числа компонент и начальной инициализации параметров, поскольку алгоритм ожидания-максимизации сходится к локальному экстремуму функции правдоподобия.

\subsection{Сравнение методов кластеризации}

Выделим следующие критерии сравнения методов кластеризации:
\begin{itemize}
	\item[---] К1 --- тип кластеризации (четкая или нечеткая);
	\item[---] К2 --- возможность обрабатывать числовые признаки;
	\item[---] К3 --- возможность обрабатывать категориальные признаки;
	\item[---] К4 --- необходимость задания числа кластеров;
	\item[---] К5 --- форма кластеров;
	\item[---] К6 --- входные данные;
	\item[---] К7 --- выходные данные;
\end{itemize}

В таблицах~\ref{tbl::clustering_compare_1} ---~\ref{tbl::clustering_compare_2} представлено сравнение рассмотренных методов кластеризации по выделенным критериям.

\begin{table}[H]
	\centering
	\caption{Сравнение методов кластеризации, часть 1}
	\begin{tabular}{|p{7cm}|l|c|c|c|l|}
		\hline
		\textbf{Метод} & \textbf{К1} & \textbf{К2} & \textbf{К3} & \textbf{К4} & \textbf{К5} \\
		\hline
		K-средних
		& Четкая
		& +
		& ---
		& +
		& Сферическая
		\\
		\hline
		K-мод
		& Четкая
		& ---
		& +
		& +
		& Сферическая
		\\
		\hline
		K-прототипов
		& Четкая
		& +
		& +
		& +
		& Сферическая
		\\
		\hline
		K-медоидов
		& Четкая
		& +
		& +
		& +
		& Произвольная
		\\
		\hline
		C-средних
		& Нечеткая
		& +
		& ---
		& +
		& Сферическая
		\\
		\hline
		DBSCAN
		& Четкая
		& +
		& ---
		& ---
		& Произвольная
		\\
		\hline
		Иерархическая кластеризация
		& Четкая
		& +
		& +
		& ---
		& Произвольная
		\\
		\hline
		Смешанная гауссовская модель (GMM)
		& Нечеткая
		& +
		& ---
		& +
		& Эллипсоидальная
		\\
		\hline
	\end{tabular}
	\label{tbl::clustering_compare_1}
\end{table}


\begin{table}[H]
	\centering
	\caption{Сравнение методов кластеризации, часть 2}
	\begin{tabular}{|p{4.5cm}|p{5.5cm}|p{5.5cm}|}
		\hline
		\textbf{Метод} & \textbf{К6} & \textbf{К7} \\
		\hline
		K-средних
		& Массив объектов, число кластеров
		& Кластеры с номерами объектов
		\\
		\hline
		K-мод
		& Массив объектов, число кластеров
		& Кластеры с номерами объектов
		\\
		\hline
		K-прототипов
		& Массив объектов, число кластеров
		& Кластеры с номерами объектов
		\\
		\hline
		K-медоидов
		& Массив объектов, число кластеров, мера расстояний
		& Кластеры с номерами объектов
		\\
		\hline
		C-средних
		& Массив объектов, число кластеров
		& Центры кластеров, матрица степеней принадлежности
		\\
		\hline
		DBSCAN
		& Массив объектов, мера расстояний, $\varepsilon$, $minPts$
		& Кластеры с номерами объектов и точки, которые не вошли ни в один кластер
		\\
		\hline
		Иерархическая кластеризация
		& Массив объектов, мера расстояний
		& Дендрограмма
		\\
		\hline
		Смешанная гауссовская модель (GMM)
		& Массив объектов, число компонент
		& Параметры компонент, матрица вероятностей принадлежности
		\\
		\hline
	\end{tabular}
	\label{tbl::clustering_compare_2}
\end{table}



\section{Методы оценки качества кластеризации}

Методы оценки качества кластеризации позволяют количественно оценить результаты кластеризации.
Принято выделять две группы методов качества кластеризации:
\begin{itemize}
	\item[---] внешние меры основаны на сравнении результата кластеризации с априори известным разделением на кластеры;
	\item[---] внутренние меры отображают качество кластеризации только по информации в данных.
\end{itemize}

Для нечетких методов кластеризации результаты можно привести к четкому разбиению, присвоив каждому объекту тот кластер, для которого степень принадлежности или вероятность максимальна.

Так как при решении задачи кластеризации обычно нет эталонного разбиения на кластеры, наиболее часто используемыми являются внутренние меры. Рассмотрим основные внутренние меры оценки качества кластеризации для разбиения множества объектов $X = \{x_1, \dots, x_n\}$ на множество кластеров $C = \{ C_1, \dots, C_k \}$.

\subsection{Компактность кластеров}

Идея данного метода заключается в том, что чем ближе друг к другу находятся объекты внутри кластеров, тем лучше разделение. Таким образом необходимо минимизировать внутриклассовое расстояние:
\begin{equation}
	WSS = \sum_{j=1}^{k}\sum_{x_i \in C_j}\|x_i - \mu_j\|^2,
\end{equation}
где $\mu_j$ --- центроид $j$-го кластера.

Чем меньше значение $WSS$, тем выше компактность кластеров, а следовательно, тем лучше качество разбиения.

\subsection{Отделимость кластеров}
В отличие от комактности, отделимость отражает, насколько далеко друг от друга расположены центра кластеров. Чем дальше находятся друг от друга разные кластеры, тем лучше. Поэтому необходимо максимизировать сумму квадратов отклонений:

\begin{equation}
	BSS = \sum_{j=1}^{k} |C_j| \cdot \|\mu_j - \bar{\mu}\|^2.
\end{equation}
где $\mu_j$ --- центроид $j$-го кластера, $\bar{\mu}$ --- глобальный центроид всего набора данных.

Чем больше значение $BSS$, тем выше разделимость кластеров, а следовательно, тем лучше качество разбиения.

\subsection{Индекс оценки силуэта}
Значение индекса оценки силуэта~\cite{Оценка_кластеризации} показывает насколько объект близок к своему кластеру по сравнению с другими кластерами.

Определим две величины:

\begin{itemize}
	\item[---] среднее расстояние от объекта $x_i \in C_j$ до других объектов своего кластера (компактность):
	\begin{equation}
		a(x_i) = \frac{1}{|C_j| - 1} \sum_{\substack{x \in C_j \\ x \neq x_i}} d(x_i, x).
	\end{equation}
	\item[---] минимальное среднее расстояние от объектов кластера до объектов другого кластера (отделимость):
	\begin{equation}
		b(x_i) = \min_{\substack{C_r \in C \\ C_r \neq C_j}}
		\left\{\frac{1}{|C_r|}\sum_{x \in C_r} d(x_i, x)\right\}.
	\end{equation}

\end{itemize}

«Силуэт» для объекта $x_i$:
\begin{equation}
	s(x_i)=\frac{b(x_i)-a(x_i)}{\max\{a(x_i),\, b(x_i)\}}.
\end{equation}

Оценка всей кластерной структуры достигается усреднением показателя по элементам:
\begin{equation}
	Sil(C)=\frac{1}{n}\sum_{i=1}^{n} s(x_i).
\end{equation}


Значения коэффициента силуэта лежат в диапазоне от $-1$ до $1$. Чем ближе значение $Sil(C)$ к 1, тем лучше: объект в среднем ближе к объектам своего кластера, чем к объектам ближайшего альтернативного кластера. Значения около 0 указывают на невыраженную кластерную структуру, отрицательные значения сигнализируют о потенциальной ошибке кластеризации.

\subsection{Индекс Калински-Харабаша}

Индекс Калински-Харабаша~\cite{Оценка_кластеризации} представляет собой отношение межкластерной дисперсии к внутрикластерной дисперсии с поправкой на число кластеров и объем данных. Вычисляется по формуле:

\begin{equation}
	CH(C) = \frac{n - k}{k - 1} \cdot \frac{\sum\limits_{C_i \in C} |C_i| \cdot ||\mu_i - \bar{\mu}||^2}{\sum\limits_{C_i \in C}\sum \limits_{x \in C_i} ||x - \mu_i||^2},
\end{equation}
где $\mu_i$ --- центроид $i$-го кластера, $\bar{\mu}$ --- глобальный центроид всего набора данных.

Числитель отражает удаленность центроидов кластеров от глобального центроида --- разделимость, знаменатель отражает сумму внутрикластерных расстояний до центроидов --- компактность. При улучшении кластеризации индекс возрастает.

\subsection{Индекс Дэвиса-Болдуина}

Индекс Дэвиса-Болдуина~\cite{Оценка_кластеризации} строится на попарном сравнении кластеров. Для каждого кластера ищется наиболее похожий на него другой кластер по отношению компактности к отделимости. Для индекса Дэвиса-Болдуина компактность вычисляется как расстояние от объектов кластера до их центроида по формуле~\ref{eq::s_i}, а отделимость --- как расстояние между центроидами.

\begin{equation}
	S(C_i) = \frac{1}{|C_i|} \sum_{x \in C_i}||x - \mu_i||,
	\label{eq::s_i}
\end{equation}
где $\mu_i$ --- центроид $i$-го кластера.

Тогда индекс Дэвиса-Болдуина высчитывается по формуле:
\begin{equation}
	DB(C) = \frac{1}{k} \sum_{C_i \in C}
	\max_{\substack{C_j \in C \\ C_j \neq C_i}}
	\left\{ \frac{S(C_i) + S(C_j)}{\|\mu_i - \mu_j\|} \right\}.
\end{equation}

Меньшее значение $DB(C)$ соответствует более компактным кластерам и более отчетливо разделенным центроидам, поэтому при улучшении качества кластеризации ее значение должно уменьшаться.

